% Page budget: 1.55 pages | Owns: negation, OBC derivation, additional-state scope, implementation, and true-parameter condition.
% WRITING GUIDE
% Purpose: Explain negation, derive the final state-level orthogonal-by-construction map, and delimit its interpretability guarantee.
% Start here: Decisions D-185, D-190, D-192, D-194, D-198, D-202, D-203, D-222, D-229, and D-232 in docs/decisions.md.
% Supporting material: Current state-level OBC derivations under scripts/gantry/orthogonal-by-construction/documentation and their thesis-readiness guidance.
% Mandatory papers: Read Györök's projection-regularization paper and orthogonal-by-construction paper in full before writing; keep their guarantees and assumptions separate.
% Research-plan anchor: Preserve Aspect 3's goal of protecting physical interpretability; state clearly how the final construction differs from the proposed regularisation.
% Current decisions: Reduced identifiable coordinates, data-derived reference tuples, a tangent-only one-step physical-state projection, per-epoch refresh, and no affine comparison arm.
% Choices to justify: Protected parameter coordinates, reference-set construction, projected rows, rank tolerance, exclusion of the affine offset, refresh frequency, and treatment of additional states.
% Implementation audit: Read model_augmentation/fit_systems/obc.py, scripts/gantry/gantry_dynamic/obc_gantry.py, obc_diagnostics.py, and the OBC attachment in model.py before fixing the derivation.
% Reference check: Verify the original projection method and separate its claims from the state-level, nonlinear-parameter, LPV, and additional-state extensions developed here.
% Cautions: docs/orthogonal-projection-plan.md describes an older soft-penalty route; one-step state orthogonality is not trajectory-output orthogonality or guaranteed true recovery.
% Planned figures: New geometric projection figure with a physical and additional-state partition inset; show the remaining uniqueness limitation as clearly as the protected direction.
% Section is finished when: The protected space, coordinate frame, gradient treatment, added-state scope, and true-recovery condition are explicit.
\section{Interpretability by Orthogonal Construction}\label{sec:obc}

Joint estimation of $\theta_{\mathrm{base}}$ and $\theta_{\mathrm{aug}}$ does not
fix the split between the baseline and the network
(Section~\ref{sec:aug_structure}).
This section makes the split unique by adapting the orthogonal-by-construction
(OBC) parametrisation of Gy\"or\"ok et al.~\cite{gyorok2026obc}, derived for
input-output models linear in their parameters, to the self-scheduled
state-space model \eqref{eq:aug_transition}.
Section~\ref{sec:obc_negation} identifies the directions in which the network
can replace the baseline, Section~\ref{sec:obc_map} constructs the map that
removes them, and Section~\ref{sec:obc_scope} states what the construction
guarantees and when it recovers the true parameters.

\subsection{Negation of the baseline}\label{sec:obc_negation}
The network writes $f_{\mathrm{aug}}$ additively into the physical rows of
\eqref{eq:aug_transition}, so it can produce any change that a parameter update
produces~\cite[Ex.~1]{gyorok2025l4dc}, \cite[Ex.~2]{gyorok2026obc}.
To first order, these changes are the columns of the transition sensitivity
\begin{equation}
 \Phi_{\bar\theta}(\tilde x,u)
 =\left.\frac{\partial f_{\mathrm{base}}(\theta_{\mathrm{base}},\tilde x,u)}
   {\partial\theta_{\mathrm{base}}}\right|_{\theta_{\mathrm{base}}=\bar\theta_{\mathrm{base}}}
   \in\R^{6\times10},
\label{eq:obc_sens}
\end{equation}
where $\bar\theta_{\mathrm{base}}$ is the expansion point.
A network write $-\Phi_{\bar\theta}\,\delta\theta$ cancels a parameter
displacement $\delta\theta$ to first order:
\begin{equation}
\begin{aligned}
 &f_{\mathrm{base}}(\bar\theta_{\mathrm{base}}+\delta\theta,\tilde x,u)
 +\bigl[f_{\mathrm{aug}}(\theta_{\mathrm{aug}},\hat x,u)
 -\Phi_{\bar\theta}(\tilde x,u)\,\delta\theta\bigr]\\
 &\quad\approx f_{\mathrm{base}}(\bar\theta_{\mathrm{base}},\tilde x,u)
 +f_{\mathrm{aug}}(\theta_{\mathrm{aug}},\hat x,u).
\end{aligned}
\label{eq:obc_negation}
\end{equation}
We call this negation.

Along these directions the objective \eqref{eq:aug_loss} cannot attribute the
dynamics to either component, so the estimate of $\theta_{\mathrm{base}}$ is
set by the start and the optimiser rather than by the data.
A penalty on the deviation from the initial parameters bounds this
drift~\cite[Sec.~3.4]{hoekstra2025lfr}, but it ties the estimate to initial
values that are themselves uncertain~\cite[Sec.~5.2]{gyorok2026obc}, and it is
not used here (Section~\ref{sec:setup}).

For the gantry, differentiating \eqref{eq:eom} with respect to the $j$-th
combination gives the acceleration sensitivity
\begin{equation}
 \frac{\partial\ddot q}{\partial\theta_{\mathrm{base},j}}
 =-M(Y)^{-1}\Bigl[\frac{\partial M(Y)}{\partial\theta_{\mathrm{base},j}}\,\ddot q
 +\frac{\partial C}{\partial\theta_{\mathrm{base},j}}\,\dot q
 +\frac{\partial K}{\partial\theta_{\mathrm{base},j}}\,q\Bigr],
\label{eq:obc_gantry_sens}
\end{equation}
which the RK4 step propagates into $\Phi_{\bar\theta}$.
A negating write therefore acts as an added inertial, damping or stiffness
force, and the factor $M(Y)^{-1}$ schedules each such direction with the
payload position.

The protected directions are those of the ten identifiable combinations
$\theta_{\mathrm{base}}$ of Section~\ref{sec:params}.
The raw parameters would add the four null directions of that section, which
leave the transition unchanged and make the stacked sensitivity rank deficient.
The training coordinates of Section~\ref{sec:params} map to
$\theta_{\mathrm{base}}$ with a nonsingular Jacobian, so the sensitivity with
respect to either has the same range.

\subsection{Orthogonal-by-construction transition}\label{sec:obc_map}
Gy\"or\"ok et al.\ remove the overlap without a penalty.
For a baseline $\hat y_k=\phi(\zeta_k)\theta_{\mathrm{base}}$ linear in its
parameters, with $\zeta_k$ the regression vector of past inputs and outputs, they
fit the stacked network output $F$ with the stacked regressor $\Phi$ and
subtract the fit~\cite[Eqs.~(8), (10), (11), (13)]{gyorok2026obc}:
\begin{equation}
\begin{aligned}
 \theta_{\mathrm{aux}}&=(\Phi^\top\Phi)^{-1}\Phi^\top F,\\
 \widetilde F&=F-\Phi\,\theta_{\mathrm{aux}},\\
 0&=\Phi^\top\widetilde F,\\
 \hat y_k&=\phi(\zeta_k)\theta_{\mathrm{base}}+f_{\mathrm{aug}}(\theta_{\mathrm{aug}},\zeta_k)
   -\phi(\zeta_k)\,\theta_{\mathrm{aux}}.
\end{aligned}
\label{eq:obc_gyorok}
\end{equation}
The third line holds for every $\theta_{\mathrm{aug}}$~\cite[Lemma~3]{gyorok2026obc}.
The research plan proposed the projection regularisation
of~\cite{gyorok2025l4dc} instead, which penalises the projected network output
with a weight that trades orthogonality against fit~\cite[Eq.~(13)]{gyorok2025l4dc};
the construction removes that weight.

Three properties of \eqref{eq:aug_transition} differ from the setting of
\eqref{eq:obc_gyorok}: $f_{\mathrm{base}}$ is rational in
$\theta_{\mathrm{base}}$ through $M(Y)^{-1}$, only the positions are measured,
and the model carries additional states.
For the first, the sensitivity \eqref{eq:obc_sens} takes the place of the
regressor, following the parameter expansion
of~\cite[Eq.~(15)]{gyorok2025l4dc}.
That paper also protects the offset of the
expansion~\cite[Eq.~(18)]{gyorok2025l4dc}; we do not.
The offset multiplies no estimated parameter, and because a discrete transition
is dominated by the propagation of the state itself, protecting it would remove
a broad state-like direction from the network that is unrelated to parameter
non-uniqueness.
The protected set is therefore the range of $\Phi_{\bar\theta}$, the
counterpart of the non-uniqueness in~\cite[Ex.~2]{gyorok2026obc}.

$\Phi_{\bar\theta}$ and $f_{\mathrm{aug}}$ are evaluated on a fixed set of
points, which~\cite[Sec.~3.2]{gyorok2026obc} permits.
A replay of the baseline on the recorded input drifts on the free axes $X$ and
$Y$ (Section~\ref{sec:aug_closed_loop}), so the points are built from the
measurements of every sample $i=1,\dots,N$ of the training records:
\begin{equation}
\begin{aligned}
 q_i&=P^{-\top}y_i,\qquad
 \dot q_i=\frac{-q_{i+2}+8q_{i+1}-8q_{i-1}+q_{i-2}}{12\,T_s},\\
 \tilde x_i&=\col(q_i,\dot q_i),\qquad \bar x_i=0,
\end{aligned}
\label{eq:obc_tuples}
\end{equation}
with $u_i$ the recorded input and the stencil kept within one record.
The stencil is fourth order because the damping directions of
\eqref{eq:obc_gantry_sens} multiply the velocity estimate.
Each tuple carries its measured $Y_i$, so $\Phi_{\bar\theta}$ is evaluated at
the scheduling value of the sample rather than expanded in $Y$, and the
protected range covers the payload positions of the training data.
The recorded input contains the feedback action, so the controller does not
enter the tuples.
The tuples are normalised as the signals of Section~\ref{sec:aug_structure}.

On these tuples, \eqref{eq:obc_gyorok} becomes
\begin{equation}
\begin{aligned}
 \Phi_{\bar\theta}(\tilde X,U)
 &=\col\bigl(\Phi_{\bar\theta}(\tilde x_1,u_1),\dots,
   \Phi_{\bar\theta}(\tilde x_N,u_N)\bigr)\in\R^{6N\times10},\\
 F_{\mathrm{aug}}
 &=\col\bigl(f_{\mathrm{aug}}(\theta_{\mathrm{aug}},\hat x_1,u_1),\dots,
   f_{\mathrm{aug}}(\theta_{\mathrm{aug}},\hat x_N,u_N)\bigr),\\
 \theta_{\mathrm{aux}}&=\Phi_{\bar\theta}(\tilde X,U)^{+}F_{\mathrm{aug}},\\
 0&=\Phi_{\bar\theta}(\tilde X,U)^\top
   \bigl(F_{\mathrm{aug}}-\Phi_{\bar\theta}(\tilde X,U)\,\theta_{\mathrm{aux}}\bigr),\\
 0&=\Phi_{\bar\theta}(\tilde X,U)^\top\,
   \frac{\partial\bigl(F_{\mathrm{aug}}-\Phi_{\bar\theta}(\tilde X,U)\,
   \theta_{\mathrm{aux}}\bigr)}{\partial\theta_{\mathrm{aug}}},
\end{aligned}
\label{eq:obc_ours}
\end{equation}
where $\hat x_i=\col(\tilde x_i,\bar x_i)$ and $^{+}$ is the Moore-Penrose
pseudoinverse.
The stack must have full column rank, so that $\theta_{\mathrm{aux}}$ is unique
and all ten directions are separated; its numerical rank is decided at the
standard tolerance relative to the largest singular value and checked at every
rebuild.
The orthogonality holds for the stack, not per tuple: one tuple gives six rows
against ten columns, and where $\Phi_{\bar\theta}(\tilde x_i,u_i)$ has full row
rank a per-tuple projection would remove the entire learned write.
The last line follows from the fourth because the stack does not depend on
$\theta_{\mathrm{aug}}$ and $\theta_{\mathrm{aux}}$ is differentiated through,
so the gradient that reaches the network is projected as well.
The stack itself is detached: $\theta_{\mathrm{base}}$ receives its gradient
only through $f_{\mathrm{base}}$ in the rollout.

The coefficient $\theta_{\mathrm{aux}}$ is computed once per objective
evaluation from the fixed tuples and reused at every step of every rollout.
Fitting it on the rollout states instead would make the correction at step $k$
depend on later states of the same rollout.
The expansion point and the stack are rebuilt at the current
$\theta_{\mathrm{base}}$ at the start of every epoch and held fixed within it.
Gy\"or\"ok et al.\ either update the expansion point at every objective
evaluation or fix it at the nominal parameters~\cite[Sec.~4]{gyorok2025l4dc}; the
first rebuilds the full stack at every evaluation, and the second does not apply
because training starts from detuned parameters (Section~\ref{sec:setup}).
Validation and prediction use the coefficient of the current network.

The corrected augmented model is
\begin{equation}
\begin{aligned}
 \begin{bmatrix}\tilde x_{k+1}\\ \bar x_{k+1}\end{bmatrix}
 &=\begin{bmatrix}f_{\mathrm{base}}(\theta_{\mathrm{base}},\tilde x_k,\hat u_k)\\ 0\end{bmatrix}\\
 &\quad+\begin{bmatrix}f_{\mathrm{aug}}(\theta_{\mathrm{aug}},\hat x_k,\hat u_k)
   -\Phi_{\bar\theta}(\tilde x_k,\hat u_k)\,\theta_{\mathrm{aux}}\\
   g_{\mathrm{aug}}(\theta_{\mathrm{aug}},\hat x_k,\hat u_k)\end{bmatrix},
\end{aligned}
\label{eq:obc_model}
\end{equation}
which replaces \eqref{eq:aug_transition} in the closed-loop rollout.
At a rollout point, the correction is the fitted combination of sensitivities
evaluated at that point, with the expansion point of the epoch, as
in~\cite[Eq.~(13)]{gyorok2026obc}.
The product $\Phi_{\bar\theta}\theta_{\mathrm{aux}}$ is the directional
derivative of the RK4 step, computed exactly by propagating a forward
sensitivity through its stages.
The baseline has zero additional-state rows, so $g_{\mathrm{aug}}$ is not
corrected.
% Planned figure (new): geometric decomposition of F_aug into its part in range(Phi) and its
% orthogonal part, on the reference stack at the expansion point, with a curved surface marking
% that the tangent is local. Inset: state partition with the corrected physical rows, the
% uncorrected additional-state rows, and their delayed influence through the recursion.

\subsection{Scope and condition for true parameter recovery}\label{sec:obc_scope}
The construction guarantees the fourth line of \eqref{eq:obc_ours} and nothing
more.
The identity is local at $\bar\theta_{\mathrm{base}}$, holds for the one-step
stack over the tuples, and uses the Euclidean inner product in normalised state
coordinates, so the state normalisation defines its metric.
The tuples lie on $\bar x=0$, so the coefficient does not see directions in
which $f_{\mathrm{aug}}$ acts only through nonzero additional states, and the
uncorrected $g_{\mathrm{aug}}$ reaches the physical rows one step later
through $f_{\mathrm{aug}}$.
The construction therefore does not make the trajectory or the output of the
augmentation orthogonal to the baseline.
The consistency and zero-covariance results
of~\cite[Thms.~16, 18]{gyorok2026obc} assume that the prediction-error criterion
is evaluated on the orthogonalised stack; the closed-loop multistep criterion
\eqref{eq:aug_loss} is not, so they do not transfer.

A unique split is not yet the true split.
Let $\theta^*_{\mathrm{base}}$ be the true parameters and $\Delta^*$ the true
one-step discrepancy on the tuples, in the same normalised coordinates,
\begin{equation}
 \Delta^*=\col\bigl(\tilde x_{i+1}-f_{\mathrm{base}}(\theta^*_{\mathrm{base}},
 \tilde x_i,u_i)\bigr)_{i=1}^{N},
\label{eq:obc_delta}
\end{equation}
where $\tilde x_{i+1}$ is the next tuple of the same record.
Suppose the data are noise-free, $f_{\mathrm{base}}$ is affine in
$\theta_{\mathrm{base}}$ between $\hat\theta_{\mathrm{base}}$ and
$\theta^*_{\mathrm{base}}$, the stack $\Phi$ of \eqref{eq:obc_ours} is evaluated
at $\theta^*_{\mathrm{base}}$, and the corrected model reproduces every
successor exactly~\cite[Assumption~6]{gyorok2026obc}.
Multiplying this one-step fit by $\Phi^{+}$ and using the fourth line of
\eqref{eq:obc_ours} gives
\begin{equation}
 \hat\theta_{\mathrm{base}}-\theta^*_{\mathrm{base}}=\Phi^{+}\Delta^*,
 \qquad
 \hat\theta_{\mathrm{base}}=\theta^*_{\mathrm{base}}
 \;\Longleftrightarrow\;
 \Phi^\top\Delta^*=0,
\label{eq:obc_bias}
\end{equation}
the state-level counterpart of~\cite[Cond.~4, Eq.~(21)]{gyorok2026obc}.
If the condition fails, the estimate is still unique but biased.
Because the training objective is the closed-loop multistep error,
\eqref{eq:obc_bias} is a one-step surrogate for the bias of the trained model,
not its value.
For a missing generalised force the condition cannot hold at a single tuple,
so it requires cancellation over the stack, which in general needs designed
excitation~\cite[Remark~5]{gyorok2026obc}.

The condition is checked with an overlap measure.
For a stacked vector $v$ over the tuples and a basis $B$,
\begin{equation}
 \rho(B;v)=\frac{\norm{BB^{+}v}}{\norm{v}},
 \qquad
 B\in\bigl\{\Phi,\;[\Phi\;\;F_{\mathrm{base}}]\bigr\},
\label{eq:obc_overlap}
\end{equation}
where $F_{\mathrm{base}}=\col\bigl(f_{\mathrm{base}}(\bar\theta_{\mathrm{base}},
\tilde x_i,u_i)\bigr)_{i=1}^{N}$ adds the unprotected offset direction.
The value $\rho^*=\rho(\Phi;\Delta^*)$, with $\Phi$ at
$\theta^*_{\mathrm{base}}$, is zero exactly when the recovery condition holds;
it needs $\theta^*_{\mathrm{base}}$ and is available in simulation only.
With $B=[\Phi\;\;F_{\mathrm{base}}]$, the same measure gives the condition
under the stricter definition that also protects the offset, and
$\rho(\Phi;F_{\mathrm{aug}})$ is the share of the network write that the
construction removes.
Section~\ref{sec:results} tests whether the split is independent of the start
and, where $\rho^*$ is small, whether $\theta^*_{\mathrm{base}}$ is recovered.
